\section{Types of discontinuities}

When a function is not continuous at a point, the failure can happen in different ways. Classifying the discontinuity helps us understand what kind of break appears in the graph and whether it can be repaired.

\subsection*{Removable discontinuity}

A removable discontinuity occurs when the limit exists, but the value at the point is missing or does not agree with the limit.

\begin{examplebox}
Let
\[
f(x)=\frac{x^2-4}{x-2},
\qquad x\ne2.
\]
For \(x\ne2\),
\[
f(x)=x+2.
\]
Therefore
\[
\lim_{x\to2}f(x)=4.
\]
The function is not defined at \(2\), but if we define \(f(2)=4\), the function becomes continuous there. The discontinuity is removable.
\end{examplebox}

Graphically, a removable discontinuity is often seen as a hole.

\subsection*{Jump discontinuity}

A jump discontinuity occurs when the left-hand and right-hand limits exist but are different.

\begin{examplebox}
Let
\[
f(x)=
\begin{cases}
0, & x<1,\\
2, & x\ge1.
\end{cases}
\]
Then
\[
\lim_{x\to1^-}f(x)=0,
\qquad
\lim_{x\to1^+}f(x)=2.
\]
The one-sided limits are different, so the two-sided limit does not exist. The function has a jump discontinuity at \(1\).
\end{examplebox}

A jump cannot be repaired by changing only the value at the point. The two sides approach different heights.

\subsection*{Infinite discontinuity}

An infinite discontinuity occurs when the function grows without bound near a point. This often happens near vertical asymptotes.

\begin{examplebox}
Consider
\[
f(x)=\frac1{x-2}.
\]
As \(x\to2^+\), the denominator is small and positive, so
\[
f(x)	o\infty.
\]
As \(x\to2^-\), the denominator is small and negative, so
\[
f(x)	o-\infty.
\]
The function has an infinite discontinuity at \(2\), and the line \(x=2\) is a vertical asymptote.
\end{examplebox}

This is different from a removable discontinuity. There is no finite value that can be assigned at \(x=2\) to make the function continuous.

\subsection*{Oscillatory behavior}

Some functions fail to have a limit because they oscillate more and more rapidly near a point.

\begin{examplebox}
The function
\[
f(x)=\sin\frac1x,
\qquad x\ne0,
\]
oscillates infinitely many times as \(x\to0\). It does not approach one number, so
\[
\lim_{x\to0}\sin\frac1x
\]
does not exist.
\end{examplebox}

This kind of example is more advanced than a simple jump, but it teaches an important lesson: a limit requires approaching one value.

\begin{summarybox}
Common types of discontinuities:
\begin{itemize}
\item Removable: the limit exists, but the value is missing or incorrect.
\item Jump: the left-hand and right-hand limits are different.
\item Infinite: the function grows without bound near the point.
\item Oscillatory: the function fails to approach one value because of persistent oscillation.
\end{itemize}
\end{summarybox}

A discontinuity is not merely a place where a formula looks difficult. It is a specific kind of failure of continuity.